%% kapa: exam=2 round=31 year=2020 period=3 subject=compensation_law subject_ko=감정평가 및 보상법규 seq=1 number=1 points=40
%% source: Q-Net 감정평가사 기출문제 (cst00309) artlSeq=5207673 file=raw/1008/5207673/3교시_감정평가 및 보상법규.pdf
%% conversion: kordoc 4.24.51 -> tools/gen2.py (auto); review=proofread-gpt-astra; flags=none
\begin{longquestion}{【문제 1】}
A시 시장 甲은 1990년에 「자연공원법」에 의하여 A시내 산지 일대 5 km\textsuperscript{2}를 ‘X시립공원’으로 지정ㆍ고시한 다음, 1992년에 X시립공원 구역을 구분하여 용도지구를 지정하는 내용의 ‘X시립공원 기본계획’을 결정ㆍ공고하였다. 甲은 2017년에 X시립공원 구역 내 10,000 m\textsuperscript{2} 부분에 다목적광장 및 휴양관 (이하 ‘이 사건 시설’이라 한다)을 설치하는 내용의 ‘X시립공원 공원계획’을 결정ㆍ고시한 다음, 2018년에 甲이 사업시행자가 되어 이 사건 시설에 잔디광장, 휴양관, 도로, 주차장을 설치하는 내용의 ‘X시립공원 공원사업’ (이하 ‘이 사건 시설 조성사업’이라 한다) 시행계획을 결정ㆍ고시하였다. 甲은 이 사건 시설 조성사업의 시행을 위하여 그 사업구역 내에 위치한 토지(이하 ‘이 사건 B토지’라 한다)를 소유한 乙과 손실보상에 관한 협의를 진행하였으나 협의가 성립하지 않자 수용재결을 신청하였다. 관할 지방 토지수용위원회의 수용재결 및 중앙토지수용위원회의 이의재결에서 모두 이 사건 B토지의 손실보상금은 1990년의 X시립공원 지정 및 1992년의 X시립공원 용도지구 지정에 따른 계획제한을 받는 상태대로 감정평가한 금액을 기초로 산정되었다. 다음 물음에 답하시오. (40점) \par
\medskip\noindent 물음 1) 乙은 위 중앙토지수용위원회의 이의재결이 감정평가에 관한 법리를 오해함으로써 잘못된 내용의 재결을 한 경우에 해당한다고 판단하고 있다. 乙이 「공익사업을 위한 토지 등의 취득 및 보상에 관한 법률」에 따라 제기할 수 있는 소송의 의의와 그 특수성을 설명하시오. (15점) \par
\medskip\noindent 물음 2) 乙이 물음 1)에서 제기한 소송에서 이 사건 B토지에 대한 보상평가는 1990년의 X시립공원 지정ㆍ고시 이전을 기준으로 하여야 한다고 주장한다. 乙의 주장은 타당한가? (10점) \par
\medskip\noindent 물음 3) 한편, 丙이 소유하고 있는 토지(이하 ‘이 사건 C토지’라 한다)는 「문화재보호법」상 보호구역으로 지정된 토지로서 이 사건 시설 조성사업의 시행을 위한 사업구역 내에 위치하고 있다. 甲은 공물인 이 사건 C토지를 이 사건 시설 조성사업의 시행을 위하여 수용할 수 있는가? (15점) \par
\begin{framedtext}
참조 조문 \par
<공익사업을 위한 토지 등의 취득 및 보상에 관한 법률> \par
제19조(토지등의 수용 또는 사용) ① <생략> \par
② 공익사업에 수용되거나 사용되고 있는 토지등은 특별히 필요한 경우가 아니면 다른 공익사업을 위하여 수용하거나 사용할 수 없다. \par
<공익사업을 위한 토지 등의 취득 및 보상에 관한 법률 시행규칙> \par
제23조(공법상 제한을 받는 토지의 평가) ① 공법상 제한을 받는 토지에 대하여는 제한받는 상태대로 평가한다. 다만, 그 공법상 제한이 당해 공익사업의 시행을 직접 목적으로 하여 가하여진 경우에는 제한이 없는 상태를 상정하여 평가한다. \par
② 당해 공익사업의 시행을 직접 목적으로 하여 용도지역 또는 용도지구 등이 변경된 토지에 대하여는 변경되기 전의 용도지역 또는 용도지구 등을 기준으로 평가한다. \par
<자연공원법> \par
제19조(공원사업의 시행 및 공원시설의 관리) ① <생략> \par
② 공원관리청은 공원사업을 하려는 경우에는 환경부령으로 정하는 기준에 따라 공원사업 시행계획을 결정하고 고시하여야 한다. \par
제22조(토지 등의 수용) ① 공원관리청은 공원사업을 하기 위하여 필요한 경우에는 공원사업에 포함되는 토지와 그 토지에 정착된 물건에 대한 소유권 또는 그 밖의 권리를 수용하거나 사용할 수 있다. \par
② 제19조제2항에 따라 공원사업 시행계획을 결정ㆍ고시한 때에는 「공익사업을 위한 토지 등의 취득 및 보상에 관한 법률」 제20조제1항 및 제22조에 따른 사업인정 및 사업인정의 고시를 한 것으로 보며, 재결 신청은 같은 법 제23조제1항 및 제28조제1항에도 불구하고 공원사업 시행계획에서 정하는 사업기간 내에 할 수 있다. \par
③ <생략> \par
④ 제1항에 따른 수용 또는 사용에 관하여는 이 법에 특별한 규정이 있는 경우를 제외하고는 「공익사업을 위한 토지 등의 취득 및 보상에 관한 법률」을 준용한다. \par
<문화재보호법> \par
제83조(토지의 수용 또는 사용) ① 문화재청장이나 지방자치단체의 장은 문화재의 보존ㆍ관리를 위하여 필요하면 지정문화재나 그 보호구역에 있는 토지, 건물, 입목(立木), 죽(竹), 그 밖의 공작물을 「공익사업을 위한 토지 등의 취득 및 보상에 관한 법률」에 따라 수용(收用)하거나 사용할 수 있다. \par
\end{framedtext}
\end{longquestion}
